% NOTE: avoid commas in \chapter titles -- a TeX Live 2025 tagging bug
% misparses them as a key list; use a colon or dash instead.
\chapter{A Sample Chapter: Sheaves and Sheafification}
\label{chap:sample}

This chapter is a demonstration. It is a short, self-contained
introduction to sheaves, in the spirit of \textcite{vakil2017}, and at
the same time it applies every building block of
Chapter~\ref{chap:blocks} in context: definitions, propositions, and
theorems with proofs; displayed equations; three figures with alt text,
one of them a commutative diagram; a table with a header row; a code
listing; links; footnotes; a block quotation; all three kinds of list;
and headings down to level four. Replace it with your own writing; keep
the patterns.

\section{Presheaves}
\label{sec:presheaves}

A sheaf is a bookkeeping device for local information: it records the
data that live over each open set of a space, and how those data
restrict to smaller open sets~\autocite{vakil2017,hartshorne1977}.

\begin{definition}[Presheaf]
  \label{def:presheaf}
  Let $X$ be a topological space. A \emph{presheaf} $\mathcal{F}$ of
  abelian groups on $X$ assigns to each open set $U \subseteq X$ an
  abelian group $\mathcal{F}(U)$, and to each inclusion of open sets
  $V \subseteq U$ a \emph{restriction} homomorphism
  $\operatorname{res}_{U,V} \colon \mathcal{F}(U) \to \mathcal{F}(V)$,
  such that $\operatorname{res}_{U,U} = \operatorname{id}$ and, for
  $W \subseteq V \subseteq U$,
  \begin{equation}
    \label{eq:functorial}
    \operatorname{res}_{V,W} \circ \operatorname{res}_{U,V}
    = \operatorname{res}_{U,W}.
  \end{equation}
\end{definition}

Elements of $\mathcal{F}(U)$ are called \emph{sections} of
$\mathcal{F}$ over $U$, and one writes $s|_{V}$ for
$\operatorname{res}_{U,V}(s)$.

\begin{example}
  \label{ex:cont}
  On any space $X$, let $\mathcal{C}(U)$ be the continuous functions
  $U \to \mathbb{R}$, with the usual restriction of functions.
  Restricting in two steps is the same as restricting in one, which is
  exactly~\eqref{eq:functorial}, so $\mathcal{C}$ is a presheaf.
\end{example}

\section{Sheaves}
\label{sec:sheaves-def}

A sheaf is a presheaf whose sections are determined by local data.

\begin{definition}[Sheaf]
  \label{def:sheaf}
  A presheaf $\mathcal{F}$ on $X$ is a \emph{sheaf} if, for every open
  $U \subseteq X$ and every open cover $U = \bigcup_i U_i$:
  \begin{enumerate}
    \item \textbf{Identity.} If $s, t \in \mathcal{F}(U)$ satisfy
          $s|_{U_i} = t|_{U_i}$ for all $i$, then $s = t$.
    \item \textbf{Gluing.} If sections $s_i \in \mathcal{F}(U_i)$ agree
          on overlaps---$s_i|_{U_i \cap U_j} = s_j|_{U_i \cap U_j}$ for
          all $i, j$---then some $s \in \mathcal{F}(U)$ has
          $s|_{U_i} = s_i$ for all $i$.
  \end{enumerate}
\end{definition}

Together the axioms say: compatible local data patch to a unique global
datum.

\begin{example}
  \label{ex:contsheaf}
  The presheaf $\mathcal{C}$ of continuous functions is a sheaf.
  Continuity is a local condition, so functions that agree on overlaps
  glue to a unique continuous function.
\end{example}

\begin{example}[A presheaf that is not a sheaf]
  \label{ex:bounded}
  On $X = \mathbb{R}$, let $\mathcal{B}(U)$ be the \emph{bounded}
  continuous functions on $U$. Cover $\mathbb{R}$ by the intervals
  $U_n = (-n, n)$ and take $s_n = x|_{U_n}$, bounded on each $U_n$. The
  $s_n$ agree on overlaps, but the only function they could glue to is
  $x$ itself, which is not bounded on $\mathbb{R}$. Gluing fails,
  because boundedness is not a local condition.
\end{example}

\begin{example}[The constant presheaf]
  \label{ex:constpre}
  Fix an abelian group $A$ with at least two elements. The
  \emph{constant presheaf} sets $\mathcal{F}(U) = A$ for every nonempty
  open $U$ and $\mathcal{F}(\emptyset) = 0$, with identity restrictions
  between nonempty sets. Suppose $U = U_1 \sqcup U_2$ is a disjoint
  union of nonempty open sets, and pick $a \neq b$ in $A$. The sections
  $a \in \mathcal{F}(U_1)$ and $b \in \mathcal{F}(U_2)$ agree on the
  overlap $U_1 \cap U_2 = \emptyset$, where both restrict to $0$---but
  no single element of $A$ restricts to $a$ on $U_1$ and to $b$ on
  $U_2$. Gluing fails.
\end{example}

Table~\ref{tab:sheaves} collects these examples and two more. The
pattern: a presheaf is a sheaf exactly when its defining condition is
local. Two standard places to read further are
\hrefu{https://stacks.math.columbia.edu/tag/006A}{the Stacks Project,
Tag 006A}~\autocite{stacksproject} and the sheaf chapter of
\hrefu{https://math.stanford.edu/~vakil/216blog/}{\emph{The Rising
Sea}}~\autocite{vakil2017}.

\begin{table}[htbp]
  \centering
  \caption[Presheaves and sheaves]{Some presheaves on
    $X = \mathbb{R}$, and whether they are sheaves.}
  \label{tab:sheaves}
  \tagpdfsetup{table/header-rows={1}}
  \begin{tabular}{lcl}
    \toprule
    \auhead{Presheaf $\mathcal{F}(U)$} & \auhead{Sheaf?} & \auhead{Reason} \\
    \midrule
    continuous functions       & yes & continuity is local \\
    smooth functions           & yes & smoothness is local \\
    locally constant functions & yes & the condition is local \\
    bounded functions          & no  & gluing fails (Ex.~\ref{ex:bounded}) \\
    constant presheaf $A$      & no  & gluing fails (Ex.~\ref{ex:constpre}) \\
    \bottomrule
  \end{tabular}
  \tablesummary{Five presheaves on the real line, each marked as a sheaf or
    not, with the reason gluing does or does not hold.}
\end{table}

\section{Stalks and germs}
\label{sec:stalks}

What does a sheaf know about a single point? Everything that happens
arbitrarily close to it.

\begin{definition}[Stalk]
  \label{def:stalk}
  The \emph{stalk} of a presheaf $\mathcal{F}$ at a point $p \in X$ is
  the direct limit over the open neighbourhoods of $p$,
  \begin{equation}
    \label{eq:stalk}
    \mathcal{F}_p \;=\; \operatorname*{colim}_{p \,\in\, U} \mathcal{F}(U),
  \end{equation}
  taken along the restriction maps (Figure~\ref{fig:stalk}).
\end{definition}

Concretely, an element of $\mathcal{F}_p$ is represented by a pair
$(U, s)$ with $p \in U$ and $s \in \mathcal{F}(U)$; two pairs $(U, s)$
and $(V, t)$ represent the same element when
\begin{equation}
  \label{eq:germ}
  s|_W = t|_W
  \qquad\text{for some open } W \subseteq U \cap V \text{ with } p \in W.
\end{equation}
This common element is the \emph{germ} of $s$ at $p$. Two sections have
the same germ exactly when they agree on some neighbourhood of $p$,
however small (Figure~\ref{fig:germ}).

\begin{figure}[tbp]
  \centering
  \includegraphics[width=0.74\linewidth,
    alt={An ellipse labelled X contains three nested dashed ellipses
      labelled U one, U two, U three, shrinking toward a marked point p.
      Below the picture, a chain of maps: F of U one, to F of U two, to
      F of U three, and so on, ending in the stalk F sub p.}]%
    {figures/stalk.pdf}
  \caption[The stalk at a point]{The stalk $\mathcal{F}_p$ is the limit
    of the groups $\mathcal{F}(U)$ over smaller and smaller
    neighbourhoods $U$ of $p$.}
  \label{fig:stalk}
\end{figure}

\begin{example}
  \label{ex:germs}
  For real-analytic functions, the germ at $p$ carries the same
  information as the Taylor series at $p$. For smooth functions the
  germ remembers strictly more.\footnote{The classic example is
  $f(x) = e^{-1/x^{2}}$, with $f(0) = 0$: it is smooth and nonzero
  arbitrarily close to $0$, yet every term of its Taylor series at $0$
  vanishes---the germ sees what the series cannot.}
\end{example}

\begin{figure}[htbp]
  \centering
  \includegraphics[width=0.74\linewidth,
    alt={Graphs of two functions, f drawn solid and g drawn dashed, over
      an interval. Inside a shaded band around a marked point p the two
      graphs coincide; outside the band they separate.}]%
    {figures/germ.pdf}
  \caption[Two functions with the same germ]{Two functions with the same
    germ at $p$: they agree on a neighbourhood of $p$, and may do as
    they please elsewhere.}
  \label{fig:germ}
\end{figure}

For a sheaf---not a mere presheaf---the germs collectively remember
everything:

\begin{proposition}[Sections are determined by their germs]
  \label{prop:germs-determine}
  Let $\mathcal{F}$ be a sheaf and $U$ open. The map
  $\mathcal{F}(U) \to \prod_{p \in U} \mathcal{F}_p$ sending a section
  to its family of germs is injective.
\end{proposition}

\begin{proof}
  If $s, t \in \mathcal{F}(U)$ have the same germ at every $p \in U$,
  then each point has a neighbourhood on which $s$ and $t$ agree. These
  neighbourhoods cover $U$, so $s = t$ by the identity axiom.
\end{proof}

This is the precise sense in which a sheaf is local: a section is
nothing more than a compatible family of germs---which is exactly the
description that sheafification, in
Section~\ref{sec:sheafification}, will turn into a definition.

A small vocabulary, for reference:
\begin{auDescription}[5.5em]
  \auitem{Section} an element of $\mathcal{F}(U)$: one piece of data over
        the open set $U$.
  \auitem{Restriction} the map $\mathcal{F}(U) \to \mathcal{F}(V)$ that
        forgets everything outside $V$.
  \auitem{Germ} the equivalence class of a section near a point.
  \auitem{Stalk} the group of all germs at a point.
\end{auDescription}

\section{Sheafification}
\label{sec:sheafification}

Natural constructions often produce presheaves that are not
sheaves---the constant presheaf is the standard example.
Sheafification repairs this, changing as little as possible.

\subsection{The universal property}
\label{sec:univ}

\begin{theorem}[Sheafification]
  \label{thm:sheafification}
  Let $\mathcal{F}$ be a presheaf on $X$. There exist a sheaf
  $\mathcal{F}^{\mathrm{sh}}$ and a morphism
  $\operatorname{sh} \colon \mathcal{F} \to \mathcal{F}^{\mathrm{sh}}$
  such that every morphism from $\mathcal{F}$ to a sheaf factors
  uniquely through $\operatorname{sh}$. Moreover, $\operatorname{sh}$
  is an isomorphism on every stalk:
  \begin{equation}
    \label{eq:stalkiso}
    \mathcal{F}_p \xrightarrow{\;\sim\;} (\mathcal{F}^{\mathrm{sh}})_p
    \qquad\text{for all } p \in X.
  \end{equation}
\end{theorem}

\begin{proof}[Proof sketch]
  Build $\mathcal{F}^{\mathrm{sh}}$ out of compatible families of
  germs: for $U$ open, let
  \begin{equation}
    \label{eq:shconstruction}
    \mathcal{F}^{\mathrm{sh}}(U) \;=\;
    \Bigl\{\, (s_p)_{p \in U} \in \textstyle\prod_{p \in U} \mathcal{F}_p
    \;:\; \text{the family is locally a section} \,\Bigr\},
  \end{equation}
  where \enquote{locally a section} means every point of $U$ has a
  neighbourhood $V$ and a section $t \in \mathcal{F}(V)$ whose germ at
  each $q \in V$ is $s_q$. The identity axiom holds because a section
  here is a family indexed by the points of $U$; the gluing axiom holds
  because membership in~\eqref{eq:shconstruction} is a local condition.
  The map $\operatorname{sh}$ sends a section to its family of germs,
  and the universal property and~\eqref{eq:stalkiso} follow by unwinding
  the definitions. Details are in
  \textcite[Section~II.1]{hartshorne1977} or \textcite{vakil2017}.
\end{proof}

Three facts worth keeping in mind:
\begin{itemize}
  \item Sheafification changes nothing at any stalk---only how local
        data assemble into sections.
  \item If $\mathcal{F}$ is already a sheaf, $\operatorname{sh}$ is an
        isomorphism.
  \item The construction is idempotent:
        $(\mathcal{F}^{\mathrm{sh}})^{\mathrm{sh}}
        \cong \mathcal{F}^{\mathrm{sh}}$, canonically.
\end{itemize}

\subsection{Worked examples}
\label{sec:worked}

\subsubsection{The constant presheaf, repaired}
\label{sec:constant-repaired}

Let $A$ be an abelian group and $\mathcal{F}$ the constant presheaf of
Example~\ref{ex:constpre}. Every stalk is $A$, so
by~\eqref{eq:shconstruction} the sheafification is the sheaf of
\emph{locally constant} functions $U \to A$, written $\underline{A}$.
On an open subset of $\mathbb{R}$ with $n$ connected components,
\begin{equation}
  \label{eq:components}
  \underline{A}(U) \cong A^{\,n},
\end{equation}
one constant per component.

\begin{corollary}
  \label{cor:twopoints}
  Let $X = \{1, 2\}$ with the discrete topology. The constant presheaf
  $\mathbb{Z}$ has global sections $\mathbb{Z}$, while its
  sheafification has global sections
  $\underline{\mathbb{Z}}(X) \cong \mathbb{Z} \oplus \mathbb{Z}$.
  Sheafification can change global sections, even though it changes no
  stalk.
\end{corollary}

\subsubsection{How to recognise a sheafification}
\label{sec:recognise}

In practice one rarely uses~\eqref{eq:shconstruction} directly. The
working tool is a standard fact: a morphism of \emph{sheaves} that is an
isomorphism on every stalk is an isomorphism. To identify
$\mathcal{F}^{\mathrm{sh}}$:
\begin{enumerate}
  \item compute the stalks of $\mathcal{F}$;
  \item find a sheaf $\mathcal{G}$ and a morphism
        $\mathcal{F} \to \mathcal{G}$ inducing an isomorphism on every
        stalk;
  \item conclude that $\mathcal{G} \cong \mathcal{F}^{\mathrm{sh}}$:
        Theorem~\ref{thm:sheafification} factors the morphism through a
        map $\mathcal{F}^{\mathrm{sh}} \to \mathcal{G}$ of sheaves that
        is still an isomorphism on every stalk---so, by the standard
        fact, an isomorphism.
\end{enumerate}

\section{Morphisms and the snake lemma}
\label{sec:snake}

Sheaves on $X$ form a category, and a remarkably well-behaved one: the
sheaves of abelian groups form an \emph{abelian} category, so kernels,
cokernels, exact sequences, and diagram chases all make
sense~\autocite[\S 2.6]{vakil2017}.

\begin{definition}[Morphism of sheaves]
  \label{def:morphism}
  A \emph{morphism} $\varphi \colon \mathcal{F} \to \mathcal{G}$ of
  sheaves on $X$ is a homomorphism
  $\varphi_U \colon \mathcal{F}(U) \to \mathcal{G}(U)$ for each open
  $U$, commuting with the restriction maps. Taking limits, $\varphi$
  induces a map $\varphi_p \colon \mathcal{F}_p \to \mathcal{G}_p$ on
  every stalk.
\end{definition}

\begin{proposition}
  \label{prop:kernel}
  For any morphism $\varphi \colon \mathcal{F} \to \mathcal{G}$, the
  presheaf $U \mapsto \ker \varphi_U$ is a sheaf.
\end{proposition}

\begin{proof}
  Identity is inherited from $\mathcal{F}$. For gluing, compatible
  sections $s_i \in \ker \varphi_{U_i}$ glue to some
  $s \in \mathcal{F}(U)$; then $\varphi_U(s)|_{U_i} =
  \varphi_{U_i}(s|_{U_i}) = 0$ for all $i$, so $\varphi_U(s) = 0$ by
  the identity axiom in $\mathcal{G}$.
\end{proof}

The cokernel is a different story: the presheaf
$U \mapsto \operatorname{coker} \varphi_U$ can fail to be a sheaf, and
the cokernel of a morphism of sheaves is \emph{defined} as its
sheafification---the first place the construction of
Section~\ref{sec:sheafification} earns its keep. Exactness is then
measured where sheaves live:

\begin{definition}[Exact sequence]
  \label{def:exact}
  A sequence of sheaves
  $\mathcal{F} \to \mathcal{G} \to \mathcal{H}$ is \emph{exact} if the
  induced sequence of stalks
  $\mathcal{F}_p \to \mathcal{G}_p \to \mathcal{H}_p$ is exact for
  every $p \in X$.
\end{definition}

\subsection{The exponential sequence}
\label{sec:exponential}

Why must the cokernel be sheafified? Because a morphism can be
surjective on every stalk yet miss sections. The classic witness lives
on the punctured plane $X = \mathbb{C} \setminus \{0\}$. Write
$\mathcal{O}_X$ for the sheaf of holomorphic functions and
$\mathcal{O}_X^{\times}$ for the nowhere-vanishing ones, and let
$\exp$ send $f$ to $e^{f}$:
\begin{equation}
  \label{eq:exponential}
  0 \longrightarrow \underline{2\pi i\,\mathbb{Z}}
    \longrightarrow \mathcal{O}_X
    \xrightarrow{\ \exp\ } \mathcal{O}_X^{\times}
    \longrightarrow 0 .
\end{equation}
The kernel is the locally constant sheaf of
Section~\ref{sec:worked}: $e^{f} = 1$ exactly when $f$ takes values in
$2\pi i\,\mathbb{Z}$. The sequence is exact: on a small disc every
nowhere-vanishing function has a logarithm, so $\exp$ is surjective on
every stalk. But over $X$ itself the function $g(z) = z$ has no
logarithm, so $\exp_X$ is not surjective on global sections---and the
image presheaf $U \mapsto \exp \mathcal{O}_X(U)$ fails gluing in
exactly the manner of Example~\ref{ex:bounded}: cover $X$ by two slit
planes, take the restrictions of $z$, and the local logarithms refuse
to patch.

\subsection{The snake lemma}
\label{sec:snakelemma}

The reward for setting up an abelian category is that the standard
tools of homological algebra apply verbatim. The first and most famous
is the snake lemma~\autocite[\S 1.7]{vakil2017}.

\begin{theorem}[Snake lemma]
  \label{thm:snake}
  In an abelian category---sheaves of abelian groups included---suppose
  the rows of the commutative diagram in Figure~\ref{fig:snake} are
  exact. Then there is a connecting morphism $\delta$ making the
  sequence
  \begin{equation}
    \label{eq:snakeseq}
    0 \to \ker f \to \ker g \to \ker h
      \xrightarrow{\ \delta\ }
      \operatorname{coker} f \to \operatorname{coker} g
      \to \operatorname{coker} h \to 0
  \end{equation}
  exact.
\end{theorem}

\begin{proof}[Proof idea]
  A diagram chase, best performed personally rather than
  read.\footnote{The chase even has a film credit: the snake lemma is
  proved on screen in the opening scene of \emph{It's My Turn} (1980).}
  Chase an element of $\ker h$ through Figure~\ref{fig:snake}: lift it
  to $B$, push it down to $B'$, observe that it comes from a unique
  element of $A'$, and define $\delta$ to be its class in
  $\operatorname{coker} f$.
\end{proof}

\begin{figure}[htbp]
  \centering
  % The diagram is drawn in LaTeX with tikz-cd (no external image). \auCDFigure
  % (defined in the class) tags it as a Figure with the Alt text below, and
  % makes tikz-cd safe under PDF math tagging. Edit the diagram here directly.
  \auCDFigure{A commutative diagram with four rows. The middle two rows are
      short exact sequences, zero to A to B to C to zero, and zero to A
      prime to B prime to C prime to zero, joined by vertical maps f, g,
      h. Above them sit the kernels of f, g, h; below them the cokernels.
      A single long arrow labelled delta runs from the kernel of h at the
      upper right, around the top and down the left of the diagram, to the
      cokernel of f at the lower left.}{%
    \begin{tikzcd}[ampersand replacement=\&, row sep=2.6em, column sep=2.2em]
    \& \ker f \arrow[r]\arrow[d] \& \ker g \arrow[r]\arrow[d] \& \ker h \arrow[d] \& \\
    0 \arrow[r] \& A  \arrow[r]\arrow[d,"f"'] \& B  \arrow[r]\arrow[d,"g"'] \& C  \arrow[r]\arrow[d,"h"'] \& 0 \\
    0 \arrow[r] \& A' \arrow[r]\arrow[d]      \& B' \arrow[r]\arrow[d]      \& C' \arrow[r]\arrow[d]      \& 0 \\
    \& \operatorname{coker} f \arrow[r] \& \operatorname{coker} g \arrow[r] \& \operatorname{coker} h \&
    \arrow[from=1-4, to=4-2, rounded corners=6pt,
      "\delta"{pos=0.5, fill=white, inner sep=2pt},
      to path={ (\tikztostart.north) -- ([yshift=0.8em]\tikztostart.north)
                -| ([xshift=-3.4em]\tikztotarget.west)
                |- (\tikztotarget.west) \tikztonodes }]
    \end{tikzcd}}
  \caption[The snake lemma]{The snake lemma. The connecting morphism
    $\delta$ gives the lemma its name.}
  \label{fig:snake}
\end{figure}

\begin{remark}
  \label{rem:cohomology}
  The snake lemma is the seed of sheaf cohomology: it is the engine
  that produces long exact sequences from short ones. Applied
  to~\eqref{eq:exponential}, it explains where the missing logarithm
  went---the failure of $\exp$ on global sections is measured by the
  first cohomology of $\underline{2\pi i\,\mathbb{Z}}$, which on the
  punctured plane records precisely the winding number.
\end{remark}

\subsubsection{By machine: a first taste of cohomology}
\label{sec:m2cohomology}

The machinery is already on your computer. In
Macaulay2~\autocite{M2}, the cohomology of line bundles on the
projective line takes three lines:

\begin{Verbatim}
R = QQ[x,y];        -- homogeneous coordinates
X = Proj R;         -- the projective line
HH^0(OO_X(2))       -- global sections: QQ^3
\end{Verbatim}

\noindent Macaulay2 answers $\mathbb{Q}^{3}$---the quadratics $x^{2}$,
$xy$, $y^{2}$---and \texttt{HH\^{}1(OO\_X(-2))} likewise returns
$\mathbb{Q}^{1}$. Where those numbers come from is the business of the
next course; that they come out of the definitions in this chapter is
the point.

As an aside on how to read this chapter---and any mathematics:
\begin{quotation}
  The substance of mathematics lies in worked examples; a definition is
  only understood once it has been turned over in the hand a few times.
\end{quotation}
